\section{AMC 10 10B 2009}
        \begin{enumerate}
        	\item (\textit{AMC 10 10B 2009}) Each morning of her five-day workweek, Jane bought either a $50$-cent muffin or a $75$-cent bagel. Her total cost for the week was a whole number of dollars. How many bagels did she buy?


 $\text{(A) } 1\qquad\text{(B) } 2\qquad\text{(C) } 3\qquad\text{(D) } 4\qquad\text{(E) } 5$


 
 \textbf{Answer: }B



\textbf{Solution 1 (Observations: Replacements)}

If Jane bought one more bagel but one fewer muffin, then her total cost for the week would increase by $25$ cents.

 If Jane bought $1$ bagel, then she bought $4$ muffins. Her total cost for the week would be $75\cdot1+50\cdot4=275$ cents, or $2.75$ dollars. Clearly, she bought one more bagel but one fewer muffin at a total cost of $3.00$ dollars. Therefore, she bought $\boxed{\textbf{(B) } 2}$ bagels.

 



\textbf{Solution 2 (Observations: Answer Choices)}


                \begin{itemize}
                    \item If Jane bought $1$ bagel, then she bought $4$ muffins. Her total cost for the week would be $75\cdot1+50\cdot4=275$ cents, or $2.75$ dollars. So, $\textbf{(A)}$ is incorrect.
                \end{itemize}


                
                \begin{itemize}
                    \item If Jane bought $2$ bagels, then she bought $3$ muffins. Her total cost for the week would be $75\cdot2+50\cdot3=300$ cents, or $3.00$ dollars. So, $\boxed{\textbf{(B) } 2}$ is correct.
                \end{itemize}


                For completeness, we will check $\textbf{(C)},\textbf{(D)},$ and $\textbf{(E)}$ too. If you decide to use this solution on the real test, then you will not need to do that, as you want to save more time.

 
                \begin{itemize}
                    \item If Jane bought $3$ bagels, then she bought $2$ muffins. Her total cost for the week would be $75\cdot3+50\cdot2=325$ cents, or $3.25$ dollars. So, $\textbf{(C)}$ is incorrect.
                \end{itemize}


                
                \begin{itemize}
                    \item If Jane bought $4$ bagels, then she bought $1$ muffin. Her total cost for the week would be $75\cdot4+50\cdot1=350$ cents, or $3.50$ dollars. So, $\textbf{(D)}$ is incorrect.
                \end{itemize}


                
                \begin{itemize}
                    \item If Jane bought $5$ bagels, then she bought $0$ muffins. Her total cost for the week would be $75\cdot5+50\cdot0=375$ cents, or $3.75$ dollars. So, $\textbf{(E)}$ is incorrect.
                \end{itemize}


                



\textbf{Solution 3 (Arithmetic)}

In this solution, all amounts are in the unit of cents.

 Note that the amount spent on muffins must end in either $00$ or $50,$ and the amount spent on bagels must end in one of $00,25,50,$ or $75.$ Furthermore, the number of muffins bought and the number of bagels bought must sum to $5.$

 We have two possible cases:

 <ol style="margin-left: 1.5em;"> \item The amounts spent on muffins and bagels both end in $00.$The number of muffins bought must be even, and the number of bagels bought must be a multiple of $4.$ 

 In this case, there are no solutions.  

 \item The amounts spent on muffins and bagels both end in $50.$The number of muffins bought must be odd, and the number of bagels bought must be $2$ more than a multiple of $4.$ 

 In this case, the only solution is $3$ muffins and $2$ bagels, from which the answer is $\boxed{\textbf{(B) } 2}.$

 </ol>



\textbf{Solution 4 (Casework, similar to others)}

We require
 \[\frac{1}{2} m + \frac{3}{4} b \in \mathbb{Z}\].We approach with standard casework. Obviously, letting $m=0,b=5$ or vice versa does not work. If $m = 1,b=4$ the right side is integral whereas the left side is not on the LHS. The same applies for the versa of that case. If $m = 2. b = 3$ then the left side is integral but the right is not of the LHS. However, when $m=3, b=2$, we have that $\frac{3}{2} + \frac{3}{2}  = 3$ is indeed integral. The answer is $\boxed{2}$.

 


	\item (\textit{AMC 10 10B 2009}) Which of the following is equal to $\dfrac{\frac{1}{3}-\frac{1}{4}}{\frac{1}{2}-\frac{1}{3}}$?


 $\text{(A) } \frac 14 \qquad \text{(B) } \frac 13 \qquad \text{(C) } \frac 12 \qquad \text{(D) } \frac 23 \qquad \text{(E) } \frac 34$


 
 \textbf{Answer: }C



\textbf{Solution 1}

Multiplying the numerator and the denominator by the same value does not change the value of the fraction.We can multiply both by $12$, getting $\dfrac{4-3}{6-4} = \boxed{\dfrac 12}$.

 Alternately, we can directly compute that the numerator is $\dfrac 1{12}$, the denominator is $\dfrac 16$, and hence their ratio is $\boxed{(C)\frac{1}{2}}$.

 



\textbf{Solution 2 (Full Solution)}

We write both the numerator and denominator with a denominator of $12$ first, since the LCM of $2,3,$ and $4$ is $3\cdot4=12$. Next, we multiply both the numerator and denominator by $12$. $\dfrac{\frac{1}{3}-\frac{1}{4}}{\frac{1}{2}-\frac{1}{3}}=\dfrac{\frac{4}{12}-\frac{3}{12}}{\frac{6}{12}-\frac{4}{12}}=\dfrac{\frac{1}{12}}{\frac{2}{12}}=\boxed{(C)\dfrac{1}{2}}$.-sosiaops

 


	\item (\textit{AMC 10 10B 2009}) Paula the painter had just enough paint for $30$ identically sized rooms.  Unfortunately, on the way to work, three cans of paint fell off her truck, so she had only enough paint for $25$ rooms.  How many cans of paint did she use for the $25$ rooms?


 $\text{(A) } 10 \qquad \text{(B) } 12 \qquad \text{(C) } 15 \qquad \text{(D) } 18 \qquad \text{(E) } 25$


 
 \textbf{Answer: }C



\textbf{Solution 1}

Losing three cans of paint corresponds to being able to paint five fewer rooms. So $\frac 35 \cdot 25 = \boxed{15}$.  The answer is $\mathrm{(C)}$.

 



\textbf{Solution 2 (proportionality)}

If $p$ is the number of cans,
 \[\frac{p}{30} = \frac{p-3}{25} \implies \frac{p}{p-3} = \frac{6}{5} \implies 5p = 6p-18 \implies p=18\].

 Thus, $p-3 = \boxed{15}$.

 


	\item (\textit{AMC 10 10B 2009}) A rectangular yard contains two flower beds in the shape of congruent isosceles right triangles. The remainder of the yard has a trapezoidal shape, as shown. The parallel sides of the trapezoid have lengths $15$ and $25$ meters. What fraction of the yard is occupied by the flower beds?


 
                    \begin{figure}[H]
                        \centering
                        \includegraphics[width=0.5\textwidth]{Problem/AMC_10_10B_2009_4_Problem_Figure_1.png}
                    \end{figure}
                    

 $\text{(A) } \frac {1}{8} \qquad \text{(B) } \frac {1}{6} \qquad \text{(C) } \frac {1}{5} \qquad \text{(D) } \frac {1}{4} \qquad \text{(E) } \frac {1}{3}$


 
 \textbf{Answer: }C



\textbf{Solution 1}

Each triangle has leg length $\frac 12 \cdot (25 - 15) = 5$ meters and area $\frac 12 \cdot 5^2 = \frac {25}{2}$ square meters.  Thus the flower beds have a total area of $25$ square meters.  The entire yard has length $25$ m and width $5$ m, so its area is $125$ square meters.  The fraction of the yard occupied by the flower beds is $\frac {25}{125} = \boxed{\frac15}$.  The answer is $\mathrm{(C)}$.

 



\textbf{Solution 2}

The length of each triangle is $\frac{25-15}{2}=5$ meters. By translating then rotating the top right triangle so that it forms a square with the top left triangle, the ratio we desire is literally the ratio of the $5*5=25$ and the total area, which is $5*25=125$. Alternatively we can use side length ratios, but each way we get $\frac{25}{125}=\frac{1}{5}$. Select $\mathrm{(C)}$.

 



\textbf{Solution 3}

Let the trapezoid be defined $ABCD$ in clockwise starting from lower left vertex and the rectangle $AEFD$ exist such that $B$ and $C$ lie on $EF$. Then, drop the altitude from $B$ and $C$ so it intersects at $G$ and $H$. Then, $GH = BC = 15$ and thus $AG = HD = 5$. The triangles are 45-45-90 and thus conjoining both we have a square of side length 5 whose area is 25, and thus the ration is $\frac{25}{25 \cdot 5} = \boxed{\frac{1}{5}}$.

 


	\item (\textit{AMC 10 10B 2009}) Twenty percent less than 60 is one-third more than what number?


 $\text{(A) } 16 \qquad \text{(B) } 30 \qquad \text{(C) } 32 \qquad \text{(D) } 36 \qquad \text{(E) } 48$


 
 \textbf{Answer: }D



\textbf{Solution}

Twenty percent less than 60 is $\frac 45 \cdot 60 = 48$. One-third more than a number \textit{n} is $\frac 43n$. Therefore $\frac 43n = 48$ and the number is $\boxed {36}$. The answer is $\mathrm{(D)}$.

 


	\item (\textit{AMC 10 10B 2009}) Kiana has two older twin brothers.  The product of their three ages is 128.  What is the sum of their three ages?


 $\text{(A) } 10 \qquad \text{(B) } 12 \qquad \text{(C) } 16 \qquad \text{(D) } 18 \qquad \text{(E) } 24$


 
 \textbf{Answer: }D



\textbf{Solution}

The age of each person is a factor of $128 = 2^7$.  So the twins could be $2^0 = 1, 2^1 = 2, 2^2 = 4, 2^3 = 8$ years of age and, consequently Kiana could be $128$, $32$, $8$, or $2$ years old, respectively.  Because Kiana is younger than her brothers, she must be $2$ years old. So the sum of their ages is $2 + 8 + 8 = \boxed{18}$.  The answer is $\mathrm{(D)}$.

 


	\item (\textit{AMC 10 10B 2009}) By inserting parentheses, it is possible to give the expression

 \[2\times3 + 4\times5\]
several values.  How many different values can be obtained?


 $\text{(A) } 2 \qquad \text{(B) } 3 \qquad \text{(C) } 4 \qquad \text{(D) } 5 \qquad \text{(E) } 6$


 
 \textbf{Answer: }C



\textbf{Solution}

The three operations can be performed on any of $3! = 6$ orders.  However, if the addition is performed either first or last, then multiplying in either order produces the same result. So at most four distinct values can be obtained.  It is easy to check that the values of the four expressions
 \begin{align*} (2\times3) + (4\times5) &= 26,\\ (2\times3 + 4)\times5 &= 50,\\ 2\times(3 + 4\times5) &= 46,\\ 2\times(3 + 4)\times5 &= 70 \end{align*}
are in fact all distinct.  So the answer is $\boxed{4}$, which is choice $\mathrm{(C)}$.

 



\textbf{Solution 2}

We use constructive counting. From $2 \times 3 + 4 \times 5$ there are four ways to place the "(". Then, there are four ways to place the ")". That accounts for 16 plausible cases.

 \textbf{Lemma:} Each value is repeated four times each. We claim the answer is 4.

 \textbf{Proof of the Lemma:}

 Very straightforward. For the single $(2) \times 3 + 4 \times 5$ obviously we see it is repeated four times. For the double, we have that $(2 \times 3) + 4 \times 5 = 2 \times 3 + (4 \times 5)$ is the same. A counting argument easily shows that if one places each of the parentheses such that it follows the Law of PEMDAS, then we obtain cases where we have a repeated sum. We handwave the rest of the proof as an interesting exercise for the reader, but the proof can be built upon PEMDAS.

 \textbf{2009 AMC 10B P7}

 The answer is $\boxed{4}$.

 


	\item (\textit{AMC 10 10B 2009}) In a certain year the price of gasoline rose by $20\%$ during January, fell by $20\%$ during February, rose by $25\%$ during March, and fell by $x\%$ during April. The price of gasoline at the end of April was the same as it had been at the beginning of January. To the nearest integer, what is $x$?


 $\text{(A) } 12 \qquad \text{(B) } 17 \qquad \text{(C) } 20 \qquad \text{(D) } 25 \qquad \text{(E) } 35$


 
 \textbf{Answer: }B



\textbf{Solution}

Let $p$ be the price at the beginning of January.  The price at the end of March was $(1.2)(0.8)(1.25)p = 1.2p.$ Because the price at the end of April was $p$, the price decreased by $0.2p$ during April, and the percent decrease was
 \[x = 100 \cdot \frac{0.2p}{1.2p} = \frac {100}{6} \approx 16.7.\]So to the nearest integer $x$ is $\boxed{17}$.  The answer is $\mathrm{(B)}$.

 



\textbf{Solution 2}

Without loss of generality, we can assume the price at the beginning of January was $\$100$. 

 When it rose by $20\%$, it became $\$120$, when it fell by $20\%$, it became $\$96$, and when it rose by $25\%$, it became $\$120$ again. 

 In order for the price at the end of April to be the same as it was at the beginning of January ($\$100$), the price must decrease by $\$20$. 

 20 is $\frac{1}{6}th$ of 120, and $\frac{1}{6} \approx 0.167 \approx 17\%$ So to the nearest integer, $x = 17$ and the answer is $\boxed{\textbf{(B) } 17}$. ~azc1027

 


	\item (\textit{AMC 10 10B 2009}) Segment $BD$ and $AE$ intersect at $C$, as shown, $AB=BC=CD=CE$, and $\angle A = \frac 52 \angle B$. What is the degree measure of $\angle D$? 


 
                    \begin{figure}[H]
                        \centering
                        \includegraphics[width=0.5\textwidth]{Problem/AMC_10_10B_2009_9_Problem_Figure_1.png}
                    \end{figure}
                    


 $\text{(A) } 52.5 \qquad \text{(B) } 55 \qquad \text{(C) } 57.5 \qquad \text{(D) } 60 \qquad \text{(E) } 62.5$


 
 \textbf{Answer: }A



\textbf{Solution}

$\triangle ABC$ is isosceles, hence $\angle ACB = \angle CAB$. 

 The sum of internal angles of $\triangle ABC$ can now be expressed as $\angle B + \frac 52 \angle B + \frac 52 \angle B = 6\angle B$, hence $\angle B = 30^\circ$, and each of the other two angles is $75^\circ$. 

 Now we know that $\angle DCE = \angle ACB = 75^\circ$.

 Finally, $\triangle CDE$ is isosceles, hence each of the two remaining angles ($\angle D$ and $\angle E$) is equal to $\frac{180^\circ - 75^\circ}2 = \frac{105^\circ}2 = \boxed{\text{(A)}\ 52.5}$.

 


	\item (\textit{AMC 10 10B 2009}) A flagpole is originally $5$ meters tall. A hurricane snaps the flagpole at a point $x$ meters above the ground so that the upper part, still attached to the stump, touches the ground $1$ meter away from the base. What is $x$?


 $\text{(A) } 2.0 \qquad \text{(B) } 2.1 \qquad \text{(C) } 2.2 \qquad \text{(D) } 2.3 \qquad \text{(E) } 2.4$


 
 \textbf{Answer: }E



\textbf{Solution 1}

The broken flagpole forms a right triangle with legs $1$ and $x$, and hypotenuse $5-x$. The \href{https://artofproblemsolving.com/wiki/index.php?title=Pythagorean\_theorem}{Pythagorean theorem} now states that $1^2 + x^2 = (5-x)^2$, hence $10x = 24$, and $x=\boxed{\text{(E)}\ 2.4}$.

 (Note that the resulting triangle is the well-known $5-12-13$ right triangle, scaled by $1/5$.)

 



\textbf{Solution 2}

A right triangle is formed with the bottom of the flagpole, the snapped part, and the ground. One leg is of length $1$ and the other is length $x$. By the \href{https://artofproblemsolving.com/wiki/index.php?title=Pythagorean\_theorem}{Pythagorean theorem}, we know that $\sqrt{x^2+1^2}$ must be the length of the snapped part of the flagpole. Observe that all the answer choices are rational. If $x$ is rational, $5-x$, which is the snapped part, must also be rational. Therefore, $1, x, 5-x$ must form a scaled Pythagorean triple. We know that $10, 24, 26$ is a Pythagorean triple, so the corresponding answer must be $1, 2.4, 2.6$. Adding together the $x$ and the snapped part, this does indeed equal $5$, so our solution is done.

 



\textbf{Solution 3}

Let $AB$ represent the flagpole in the diagram above. After the flagpole breaks at point $D$, its tip lies at point $C$. Since none of the flagpole is destroyed, we know that $DA=DC$. Therefore, triangle $\triangle ADC$ is isosceles.

 Draw the altitude $DE \perp AC$. Since $\triangle ADC$ is isosceles, we know that $AE = EC$. Also note that $\triangle AED \sim \triangle ABC$. Therefore,
 \begin{align*} AD &= AE \times \frac{AD}{AE} \\ &= \frac{AC}{2} \times \frac{AC}{AB} \\ &= \frac{AC^2}{2 AB} \\ &= \frac{AB^2 + BC^2}{2 AB} \end{align*}


 Since $AB = 5$ and $BC = 1$, we have that $AD = \frac{5^2 + 1^2}{2 \cdot 5} = 2.6$, and thus $x = AB - AD = \boxed{2.4}$.

 


	\item (\textit{AMC 10 10B 2009}) How many $7$-digit palindromes (numbers that read the same backward as forward) can be formed using the digits $2$, $2$, $3$, $3$, $5$, $5$, $5$?


 $\text{(A) } 6 \qquad \text{(B) } 12 \qquad \text{(C) } 24 \qquad \text{(D) } 36 \qquad \text{(E) } 48$


 
 \textbf{Answer: }A



\textbf{Solution}

A seven-digit palindrome is a number of the form $\overline{abcdcba}$. Clearly, $d$ must be $5$, as we have an odd number of fives. We are then left with $\{a,b,c\} = \{2,3,5\}$. There are $3!$ permutations of these three numbers, since each is reflected over the midpoint we only have to count the first there. Each of the $\boxed{(A)6}$ permutations of the set $\{2,3,5\}$ will give us one palindrome.

 



\textbf{Solution 2}

Say we have a 2 first. Then, we have a 2 pinned as the last digit, so we have to fill in the remaining digits with only 3's and 5's. We have 2 options for the second digit then, and the rest is fixed. This means that we have $2$ ways for this case.

 Say we have a 3 first. By symmetry, this is the same as the 2 cases, so we have $2$ ways.

 Say we have a 5 first. We then have a 5 in the middle. We can either have a 2 second or a 3 second. So we have $2$ ways.

 This means that our answer is $2+2+2=\boxed{\textbf{(A)}6}$

 



\textbf{Solution 3 (bijection)}

\textbf{Lemma:} 5 must be the middle digit.

 \textbf{Proof of the Lemma:}

 If it wasn't, then one 2 or 3 must be the middle digit. This then forces less 2 or 3 to reach \textit{palindromic equilibrium}, which is the number of viable digits allowed in which a palindrome can be formed. Indeed, 5 must be the middle.

 \textbf{2009 AMC 10B P11}

 Place 5 in the middle. Then consider the six digit number formed by removing 5. This integer must be a palindrome, and we have two each of 2,3,5 to make this palindrome. There are three choices for the first digit, two for the second, and one for the third, fourth, fifth, and sixth. Thus, the count is $3 \cdot 2 = \boxed{6}$.

 


	\item (\textit{AMC 10 10B 2009}) Distinct points $A$, $B$, $C$, and $D$ lie on a line, with $AB=BC=CD=1$. Points $E$ and $F$ lie on a second line, parallel to the first, with $EF=1$. A triangle with positive area has three of the six points as its vertices. How many possible values are there for the area of the triangle?


 $\text{(A) } 3 \qquad \text{(B) } 4 \qquad \text{(C) } 5 \qquad \text{(D) } 6 \qquad \text{(E) } 7$


 
 \textbf{Answer: }A



\textbf{Solution 1}

Consider the classical formula for triangle area: $\frac 12 \cdot b \cdot h$. Each of the triangles that we can make has exactly one side lying on one of the two parallel lines. If we pick this side to be the base, the height will always be the same - it will be the distance between the two lines.

 Hence each area is uniquely determined by the length of the base. And it can easily be seen, that the only possible base lengths are $1$, $2$, and $3$. Therefore there are only $\boxed{(A)3}$ possible values for the area.

 (To be more precise in the last step, the possible base lengths are $AB=BC=CD=EF=1$, $AC=BD=2$, and $AD=3$.)

 



\textbf{Solution 2}

No matter what how we draw a triangle by selecting three non-linear points, its height will always remain the same. Therefore, we will only get different areas with different base-lengths. The possibilities are $1$, $2$, and $3$ units for a total of $\boxed{\text{(A) } 3}$.

 



\textbf{Solution 3}

The orthogonal distance from two points on each of the parallel lines is equivalent to the height of any triangle formed via the points, hence its area is dependent on the base. From a segment of length 3 that is trisected by say points $A,B,C$, to any external points $E,F$, their placements do not matter. WLOG, assume it is $E$ for which the triangle is positioned at. Then, we have $\boxed{3}$ values of the base since the height is the same, particularly bases of length 1, 2, or 3.

 


	\item (\textit{AMC 10 10B 2009}) As shown below, convex pentagon $ABCDE$ has sides $AB=3$, $BC=4$, $CD=6$, $DE=3$, and $EA=7$. The pentagon is originally positioned in the plane with vertex $A$ at the origin and vertex $B$ on the positive $x$-axis. The pentagon is then rolled clockwise to the right along the $x$-axis. Which side will touch the point $x=2009$ on the $x$-axis? 


 
                    \begin{figure}[H]
                        \centering
                        \includegraphics[width=0.5\textwidth]{Problem/AMC_10_10B_2009_13_Problem_Figure_1.png}
                    \end{figure}
                    


 $\text{(A) } \overline{AB} \qquad \text{(B) } \overline{BC} \qquad \text{(C) } \overline{CD} \qquad \text{(D) } \overline{DE} \qquad \text{(E) } \overline{EA}$


 
 \textbf{Answer: }C



\textbf{Solution}

The perimeter of the polygon is $3+4+6+3+7 = 23$. Hence as we roll the polygon to the right, every $23$ units the side $\overline{AB}$ will be the bottom side.

 We have $2009 = 23 \times 87 + 8$. Thus at some point in time we will get the situation when $A=(2001,0)$ and $\overline{AB}$ is the bottom side. Obviously, at this moment $B=(2004,0)$. 

 After that, the polygon rotates around $B$ until point $C$ hits the $x$ axis at $(2008,0)$.

 And finally, the polygon rotates around $C$ until point $D$ hits the $x$ axis at $(2014,0)$.At this point the side $\boxed{\overline{CD}}$ touches the point $(2009,0)$. So the answer is $\boxed{C}$

 



\textbf{Solution 2: Mod Arithmetic}

The perimeter is $23$ and $2009\equiv8($mod $23)$, so it will end up on side $AB$ + a total of 8 more units. $4<8$, but $4+6=10>8$, so it ends on side $CD$ for an answer of  $\boxed{C}$.

 



\textbf{Solution 3 (similar to [1])}

The acquired distance per a revolution among the centre of the pentagon considers 23 units from position $A$. Thus, because $2009 \equiv 8 \pmod{23}$ and $\operatorname{quo}(2009,23) = 87$, it follows that $87 \cdot 23 = 2001$ slots are covered continuously before $A$ is positioned on the lefthand vertex once more, particularly, at $(2001,0)$. Thus, $B(2004,0)$ so in a rotation $C(2008,0)$ and thus $\boxed{CD}$ has the point $(2009,0)$.

 


	\item (\textit{AMC 10 10B 2009}) On Monday, Millie puts a quart of seeds, $25\%$ of which are millet, into a bird feeder.  On each successive day she adds another quart of the same mix of seeds without removing any seeds that are left.  Each day the birds eat only $25\%$ of the millet in the feeder, but they eat all of the other seeds.  On which day, just after Millie has placed the seeds, will the birds find that more than half the seeds in the feeder are millet?


 $\text{(A) } \text{Tuesday} \qquad \text{(B) } \text{Wednesday} \qquad \text{(C) } \text{Thursday} \qquad \text{(D) } \text{Friday} \qquad \text{(E) } \text{Saturday}$


 
 \textbf{Answer: }D



\textbf{Solution 1}

On Monday, day 1, the birds find $\frac 14$ quart of millet in the feeder.  On Tuesday they find
 \[\frac 14 + \frac 34 \cdot \frac 14\]quarts of millet.  On Wednesday, day 3, they find
 \[\frac 14 + \frac 34 \cdot \frac 14 + \left(\frac34\right)^2 \cdot \frac 14\]quarts of millet.  The number of quarts of millet they find on day $n$ is
 \[\frac 14 + \frac 34 \cdot \frac 14 + \left(\frac34\right)^2 \cdot \frac 14 + \cdots + \left(\frac 34\right)^{n-1} \cdot \frac 14 = \frac {(\frac 14)(1 - (\frac 34)^n)}{1 - \frac 34} = 1 - \left(\frac 34\right)^n .\]The birds always find $\frac 34$ quart of other seeds, so more than half the seeds are millet if $1 - \left(\frac 34\right)^n > \frac 34$, that is, when $\left(\frac 34\right)^n < \frac 14$.  Because $\left(\frac 34\right)^4 = \frac {81}{256} > \frac 14$ and $\left(\frac 34\right)^5 = \frac {243}{1024} < \frac 14$, this will first occur on day $5$ which is $\boxed {\text{Friday}}$. The answer is $\mathrm{(D)}$.

 



\textbf{Solution 2 (Brute Force)}

Notice that, every day, there is always $\frac34$ quarts of other seeds in the feeder, so we simply need to calculate when the quantity of millet seeds is greater than $\frac34$ quarts. Every day that elapses, a fourth of the millet is taken away (eaten by birds), and then a fourth of a quart is added (by Millie). Thus, if $n$ was the quantity of millet seeds the previous day, then on the next day, the quantity is $\dfrac34n+\dfrac14=\dfrac14\left(3n+1\right).$ 
Monday: $\dfrac14\quad$

Tuesday: $\dfrac14\left(3\cdot\dfrac14+1\right)=\dfrac{7}{16}$

Wednesday: $\dfrac14\left(3\cdot\dfrac{7}{16}+1\right)=\dfrac{37}{64}$

Thursday: $\dfrac14\left(3\cdot\dfrac{37}{64}+1\right)=\dfrac{175}{256}$

Friday: $\dfrac14\left(3\cdot\dfrac{175}{256}+1\right)=\dfrac{781}{1024}>\dfrac{768}{1024}=\dfrac34.$

Brute forcing, we see that it takes $4$ days. Thus, the answer is $\boxed{\text{Friday}}$, or $\boxed{\text{(D)}}$. ~ aidan0626, lucaswujc, & Technodoggo
Note: Yes there is a pattern to it, $1-\left(\frac{3}{4}\right)^x$ where $x$ is the number of days, as shown in Solution 1. This solution is only if you didn't notice the pattern. -aidan0626

 



\textbf{Solution 3}

Let $a_n$ denote the the amount of millet that the birds find on day $n$. Here $n\ge1$ and $n=1$ corresponds to Monday. Similarly, let $b_n$ denote the amount of other seeds that the birds find on day $n$. Notice that $a_n$ and $b_n$ satisfy the following:

 \begin{cases}a\_1=\frac{1}{4}\\a\_n=\frac{3}{4}a\_{n-1}+\frac{1}{4}\quad (n>1)\end{cases}

 


 
 \[b_n=\frac{3}{4}\]

 
We can solve for $a_n$ using difference equations to get $a_n=-\left(\frac{3}{4}\right)^n+1$. We want to find the least $n$ for which $a_n>\frac{1}{2}(a_n+b_n)$. In other words, $a_n>b_n$. Substituting in what we know about $a_n$ and $b_n$ we get:

 

 \[-\left(\frac{3}{4}\right)^n+1>\frac{3}{4}\]

 
 \[\left(\frac{3}{4}\right)^n<\frac{1}{4}\]

 
 \[3^n<4^{n-1}\]

 
 \[n\ge5\]

 
Hence, the answer is $\boxed{\text{Friday}}$, or $\boxed{\text{(D)}}$.

 


	\item (\textit{AMC 10 10B 2009}) When a bucket is two-thirds full of water, the bucket and water weigh $a$ kilograms. When the bucket is one-half full of water the total weight is $b$ kilograms. In terms of $a$ and $b$, what is the total weight in kilograms when the bucket is full of water?


 $\text{(A) } \frac23a + \frac13b \qquad \text{(B) } \frac32a - \frac12b \qquad \text{(C) } \frac32a + b \qquad \text{(D) } \frac32a + 2b \qquad \text{(E) } 3a - 2b$


 
 \textbf{Answer: }E



\textbf{Solution 1}

Let $x$ be the weight of the bucket and let $y$ be the weight of the water in a full bucket.  Then we are given that $x + \frac 23y = a$ and $x + \frac 12y = b$.  Hence $\frac 16y = a-b$, so $y = 6a-6b$.  Thus $x = b - \frac 12 (6a-6b) = -3a + 4b$.  Finally $x + y = \boxed {3a-2b}$. The answer is $\mathrm{(E)}$.

 



\textbf{Solution 2}

Imagine that we take three buckets of the first type, to get rid of the fraction. We will have three buckets and two buckets' worth of water.

 On the other hand, if we take two buckets of the second type, we will have two buckets and enough water to fill one bucket.

 The difference between these is exactly one bucket full of water, hence the answer is $3a-2b$.

 



\textbf{Solution 3}

We are looking for an expression of the form $xa + yb$.

 We must have $x+y=1$, as the desired result contains exactly one bucket. Also, we must have $\frac 23 x + \frac 12 y = 1$, as the desired result contains exactly one bucket of water.

 At this moment, it is easiest to check that only the options (A), (B), and (E) satisfy $x+y=1$, and out of these only (E) satisfies the second equation.

 Alternatively, we can directly solve the system, getting $x=3$ and $y=-2$.

 



\textbf{Solution 4}

Since $a$ is one bucket plus two-thirds of the total amount of water, and $b$ is one bucket plus one-half of the total amount of water, $a-b$ would equal $\cancel{\text{bucket}}+\frac23\cdot\text{water}-(\cancel{\text{bucket}}+\frac12\cdot\text{water})=\frac16\cdot\text{water}$. Therefore, $a-b$ is one-sixth of the total mass of the water. Starting from $a$, we add two one-sixths of the total amount of the water to become full. Therefore, the full bucket of water is $a+2(a-b)=a+2a-2b=3a-2b\Rightarrow\textbf{(E)}$

 -\href{https://artofproblemsolving.com/wiki/index.php?title=User:Sweetmango77}{SweetMango77}

 



\textbf{Solution 5 (Using answer choices+Very quick bash)}

WLOG let the bucket hold 6 units of water. Then b=4, a=3. Then test out all the answer choices to arrive at the answer choice (E)~PEKKA

 


	\item (\textit{AMC 10 10B 2009}) Points $A$ and $C$ lie on a circle centered at $O$, each of $\overline{BA}$ and $\overline{BC}$ are tangent to the circle, and $\triangle ABC$ is equilateral. The circle intersects $\overline{BO}$ at $D$. What is $\frac{BD}{BO}$?


 $\text{(A) } \frac {\sqrt2}{3} \qquad \text{(B) } \frac {1}{2} \qquad \text{(C) } \frac {\sqrt3}{3} \qquad \text{(D) } \frac {\sqrt2}{2} \qquad \text{(E) } \frac {\sqrt3}{2}$


 
 \textbf{Answer: }B



\textbf{Solution 1}


                    \begin{figure}[H]
                        \centering
                        \includegraphics[width=0.5\textwidth]{Solution/AMC_10_10B_2009_16_Solution_Figure_1.png}
                    \end{figure}
                    

 As $\triangle ABC$ is equilateral, we have $\angle BAC = \angle BCA = 60^\circ$, hence $\angle OAC = \angle OCA = 30^\circ$. Then $\angle AOC = 120^\circ$, and from symmetry we have $\angle AOB = \angle COB = 60^\circ$. Thus, this gives us $\angle ABO = \angle CBO = 30^\circ$.

 We know that $DO = AO$, as $D$ lies on the circle. From $\triangle ABO$ we also have $AO = BO \sin 30^\circ = \frac{BO}2$, Hence $DO = \frac{BO}2$, therefore $BD = BO - DO = \frac{BO}2$, and $\frac{BD}{BO} = \boxed{\frac 12 \Longrightarrow B}$.

 



\textbf{Solution 2}


                    \begin{figure}[H]
                        \centering
                        \includegraphics[width=0.5\textwidth]{Solution/AMC_10_10B_2009_16_Solution_Figure_2.png}
                    \end{figure}
                    

 As in the previous solution, we find out that $\angle AOB = \angle COB = 60^\circ$. Hence $\triangle AOD$ and $\triangle COD$ are both equilateral.

 We then have $\angle SCD = \angle SAD = 30^\circ$, hence $D$ is the incenter of $\triangle ABC$, and as $\triangle ABC$ is equilateral, $D$ is also its centroid. Hence $2 \cdot SD = BD$, and as $SD = SO$, we have $2\cdot SD = SD + SO = OD$, therefore $BD=OD$, and as before we conclude that $\frac{BD}{BO} = \boxed{\frac 12}$.

 



\textbf{Solution 3}

$\triangle BAO\cong\triangle BCO$ by SSS congruence, so $\angle ABO = \angle CBO = \frac{60}{2} = 30 ^\circ$. Since $BA$ is tangent to the circle, it is perpendicular to $AO$. This means that $\triangle BAO$ is a 30-60-90 triangle. The ratio of the side-lengths of a 30-60-90 triangle is $1:\sqrt3:2$, so $BO=2$. $BD = BO-DO$. Since $DO$ is the radius of the circle, $DO=1$ and $BD = 2-1=1$. Hence, $\frac{BD}{BO} = \frac{1}{2}$, and the answer is $\boxed{\textbf{(B) } \frac{1}{2}}$

 


	\item (\textit{AMC 10 10B 2009}) Five unit squares are arranged in the coordinate plane as shown, with the lower left corner at the origin. The slanted line, extending from $(c,0)$ to $(3,3)$, divides the entire region into two regions of equal area. What is $c$?

                    \begin{figure}[H]
                        \centering
                        \includegraphics[width=0.5\textwidth]{Problem/AMC_10_10B_2009_17_Problem_Figure_1.png}
                    \end{figure}
                    


 $\text{(A) } \frac 12 \qquad \text{(B) } \frac 35 \qquad \text{(C) } \frac 23 \qquad \text{(D) } \frac 34 \qquad \text{(E) } \frac 45$


 
 \textbf{Answer: }C



\textbf{Solution 1}

For $c\geq 1.5$ the shaded area is at most $1.5$, which is too little. Hence $c<1.5$, and therefore the point $(2,1)$ is indeed inside the shaded part, as shown in the picture.

 Then the area of the shaded part is one less than the area of the triangle with vertices $(c,0)$, $(3,0)$, and $(3,3)$. The area of the entire triangle is $\frac{3(3-c)}2$, therefore the area of the shaded part is $\frac{7-3c}{2}$.

 The entire figure has area $5$, hence we want the shaded part to have area $\frac 52$. Solving for $c$, we get $c=\boxed{\frac 23}$.  The answer is $\mathrm{(C)}$ ( what a coincidence ;) ).

 



\textbf{Solution 2}

The unit square is of area 1, so the five unit squares have area 5.Therefore the shaded space must occupy 2.5. The missing unit square is of area 1, and if reconstituted the original triangle would be of area 3.5.It can then be inferred:$(3-c) * 3 = 7$.

 $3-c=\frac{7}{3}$, so $3-\frac{7}{3}=c$.

 $3-\frac{7}{3} = \frac{9-7}{3} = \frac{2}{3}$. $\boxed{\mathrm{(C)}}$.

 



\textbf{Solution 3}

The shaded space of the object can become a triangle by adding a unit square to its bottom. This triangle has a base length of $3-c$ and a height of $3$. The area of this triangle region is now (using the formula $A=bh/2$ for a triangle)$(9-3c)/2$. But, remember that we have to subtract the area of the extra unit square from this area to get the area of the shaded region. 

 If you add 3 unit squares to the top of the unshaded portion, the figure is now a right trapezoid. Using the formula for a trapezoid, you can derive that its area is $(9+3c)/2$. After you subtract the area of the 3 extra unit squares, we have derived the area of the unshaded portion.

 Since these areas are both equal, we set them equal to each other and solve for $c$. $(9-3c)/2 -1 = (9+3c)/2 -3$. $c$ is now solved to be $\frac{2}{3}$. $\boxed{\mathrm{(C)}}$.

 



\textbf{Solution 4 (Brute Force)}

We are looking for the area of the shaded region to be $\frac 52$. We start by testing $(A) \frac 12$. The area of the shaded region would be $\frac{(3-\frac 12)(3) }{2}-1=\frac {11}{4}$ when $c=\frac 12$. This does not match our wanted answer.

 We try $(B) \frac 35$ next. The area of the shaded region would be $\frac{(3-\frac 35)(3) }{2}-1=\frac {13}{5}$ when $c=\frac 35$. This also does not match our desired answer.

 We then try $(C) \frac 23$. The area of the shaded region would be $\frac{(3-\frac 23)(3) }{2}-1=\frac {5}{2}$ when $c=\frac 23$. Therefore, our answer is $\boxed{\mathrm{(C) \frac 23}}$.

 



\textbf{Solution 5}

Use the shoelace theorem to get the area of the unshaded portion in terms of a. Divide this area by the whole area which is 5. This fraction should be equal to $1/2$. You can easily solve for the answer from there.-mewto

 


	\item (\textit{AMC 10 10B 2009}) Rectangle $ABCD$ has $AB=8$ and $BC=6$. Point $M$ is the midpoint of diagonal $\overline{AC}$, and $E$ is on $AB$ with $\overline{ME}\perp\overline{AC}$. What is the area of $\triangle AME$?


 $\text{(A) } \frac{65}{8} \qquad \text{(B) } \frac{25}{3} \qquad \text{(C) } 9 \qquad \text{(D) } \frac{75}{8} \qquad \text{(E) } \frac{85}{8}$


 
 \textbf{Answer: }D



\textbf{Solution 1 (Coordinate Geometry)}

Set $A$ to $(0,0)$. Since $M$ is the midpoint of the diagonal, it would be $(4,3)$. The diagonal $AC$ would be the line $y = \frac{3x}{4}$. Since $ME$ is perpendicular to $AC$, its line would be in the form $y = -\frac{4x}{3} + b$. Plugging in $4$ and $3$ for $x$ and $y$ would give $b = \frac{25}{3}$. To find the x-intercept of $y = -\frac{4x}{3} + \frac{25}{3}$ we plug in $0$ for $y$ and get $x = \frac{25}{4}$. Then, using the Shoelace Formula for $(0,0)$ , $(4,3)$, and $(\frac{25}{4}, 0)$, we find the area is $\frac{75}{8}$.

 



\textbf{Solution 2}


                    \begin{figure}[H]
                        \centering
                        \includegraphics[width=0.5\textwidth]{Solution/AMC_10_10B_2009_18_Solution_Figure_1.png}
                    \end{figure}
                    

 By the \href{https://artofproblemsolving.com/wiki/index.php?title=Pythagorean\_theorem}{Pythagorean theorem} we have $AC=10$, hence $AM=5$. 

 The triangles $AME$ and $ABC$ have the same angle at $A$ and a right angle, thus all their angles are equal, and therefore these two triangles are similar.

 The ratio of their sides is $\frac{AM}{AB} = \frac 58$, hence the ratio of their areas is $\left( \frac 58 \right)^2 = \frac{25}{64}$.

 And as the area of triangle $ABC$ is $\frac{6\cdot 8}2 = 24$, the area of triangle $AME$ is $24\cdot \frac{25}{64} = \boxed{ \frac{75}8 }$.

 



\textbf{Solution 3 (Only Pythagorean Theorem)}


                    \begin{figure}[H]
                        \centering
                        \includegraphics[width=0.5\textwidth]{Solution/AMC_10_10B_2009_18_Solution_Figure_1.png}
                    \end{figure}
                    

 Draw $EC$ as shown from the diagram. Since $AC$ is of length $10$, we have that $AM$ is of length $5$, because of the midpoint $M$. Through the Pythagorean theorem, we know that $AE^2 = AM^2 + ME^2 \implies 25 + ME^2$, which means $AE = \sqrt{25 + ME^2}$. Define $ME$ to be $x$ for the sake of clarity. We know that $EB = 8 - \sqrt{25 + x^2}$. From here, we know that $CE^2 = CB^2 + BE^2 = ME^2 + MC^2$. From here, we can write the expression $6^2 + (8 - \sqrt{25 + x^2})^2 = 5^2 + x^2 \implies 36 + (64 - 16\sqrt{25 + x^2} + 25 + x^2) = 25 + x^2 \implies 100 - 16\sqrt{25 + x^2} \implies \sqrt{25 + x^2} = \frac{25}{4} \implies 25 + x^2 = \frac{625}{16} \implies 400 + 16x^2 = 625 \implies 16x^2 = 225 \implies x = \frac{15}{4}$. Now, remember $CE \neq \frac{15}{4}$. $x = \frac{15}{4} = ME$, since we set $x = ME$ in the start of the solution. Now to find the area $\frac{15}{4} \cdot 5 \cdot \frac{1}{2} = \frac{75}{8} = CE$

 



\textbf{Solution 4 (Similarity)}

We know $AM = \frac{10}{2} = 5$ by the Pythagorean theorem, and furthermore, $\triangle AME$ is similar to $\triangle ABC$. Therefore, $ME = 5 \cdot \frac{6}{8} = \frac{15}{4}$, and the area of the triangle is $5 \cdot \frac{15}{4} \cdot \frac{1}{2} = \boxed{\frac{75}{8}}$.

 



\textbf{Solution 5 (Pythagorean Theorem)}


                    \begin{figure}[H]
                        \centering
                        \includegraphics[width=0.5\textwidth]{Solution/AMC_10_10B_2009_18_Solution_Figure_1.png}
                    \end{figure}
                    

 By the \href{https://artofproblemsolving.com/wiki/index.php?title=Pythagorean\_Theorem}{Pythagorean Theorem}, we claim that $AC = 10$. It then follows that $AM \cong MC = 5.$ 

 Because we have $AM \cong MC, \angle AME \cong \angle CME,$ and reflexive side $EM$, it follows that $\triangle AME \cong \triangle CME.$ By CPCTC, we have $AE \cong EC.$ For the sake of simplicity, we'll call those side lengths $x$. Also, since $AE = x,$ we get $BE = 8 - x.$ We can now set up the Pythagorean theorem on $\triangle EBC$: 
 \[(8 - x)^2 + 6^2 = x^2.\] Combining like terms and simplifying gives $-16x + 100 = 0$ so $x = \frac{25}{4}.$ 

 It helps to think that in order to find $[AME],$ we must have $\overline{MC}$ and $\overline{EM}.$ Let $EM = y.$ Applying the Pythagorean Theorem to $\triangle CME$ gives 
 \[5^2 + y^2 = \left(\frac{25}{4} \right)^2.\] Solving for $y$ (this is not that difficult) gives $y = \frac{15}{4}.$ So, the area of $\triangle AME$ is $\frac{\frac{15}{4} \cdot 5}{2} = \frac{75}{8} \implies \boxed{D}.$

 



\textbf{Solution 6}

The diagonal $AC = 10$ so $AM = 5$. Drop the altitude from $M$ to $AB$ call it $N$. Then, $MN = 3$, because $M$ is the centre of the rectangle. Thus, $NA = 4$. Now, consider $EMN$. Quadrilateral $EBCM$ is cyclic and thus if $\angle MEA = a$, then $\angle MEB = 180-a$ and thus $\angle ACB = a$, The two triangles are similar, and because both are 3-4-5 right triangles, $4x = 3 \implies x = \frac{3}{4}$ and thus $NE = 3 \cdot \frac{3}{4} = \frac{9}{4}$. Thus the base of $AME$ is $4 + \frac{9}{4} = \frac{25}{4}$ and the height is 3, and thus the area is $\frac{1}{2} \cdot \frac{25}{4} \cdot 3 = \boxed{\frac{75}{8}}$.

 


	\item (\textit{AMC 10 10B 2009}) A particular $12$-hour digital clock displays the hour and minute of a day. Unfortunately, whenever it is supposed to display a $1$, it mistakenly displays a $9$. For example, when it is 1:16 PM the clock incorrectly shows 9:96 PM. What fraction of the day will the clock show the correct time?


 $\text{(A) } \frac 12 \qquad \text{(B) } \frac 58 \qquad \text{(C) } \frac 34 \qquad \text{(D) } \frac 56 \qquad \text{(E) } \frac 9{10}$


 
 \textbf{Answer: }A



\textbf{Solution 1}

The clock will display the incorrect time for the entire hours of $1, 10, 11$ and $12$.  So the correct hour is displayed $\frac 23$ of the time.  The minutes will not display correctly whenever either the tens digit or the ones digit is a $1$, so the minutes that will not display correctly are $10, 11, 12, \dots, 19$ and $01, 21, 31, 41,$ and $51$.  This amounts to fifteen of the sixty possible minutes for any given hour.  Hence the fraction of the day that the clock shows the correct time is $\frac 23 \cdot \left(1 - \frac {15}{60}\right) = \frac 23 \cdot \frac 34  = \boxed{\frac 12}$.  The answer is $\mathrm{(A)}$.

 



\textbf{Solution 2}

The required fraction is the number of correct times divided by the total times. There are 60 minutes in an hour and 12 hours on a clock, so there are 720 total times.

 We count the correct times directly; let a correct time be $x:yz$, where $x$ is a number from 1 to 12 and $y$ and $z$ are digits, where $y<6$. There are 8 values of $x$ that will display the correct time: 2, 3, 4, 5, 6, 7, 8, and 9. There are five values of $y$ that will display the correct time: 0, 2, 3, 4, and 5. There are nine values of $z$ that will display the correct time: 0, 2, 3, 4, 5, 6, 7, 8, and 9. Therefore there are $8\cdot 5\cdot 9=40\cdot 9=360$ correct times.

 Therefore the required fraction is $\frac{360}{720}=\frac{1}{2}\Rightarrow \boxed{\mathrm{(A)}}$.

 


 


	\item (\textit{AMC 10 10B 2009}) Triangle $ABC$ has a right angle at $B$, $AB=1$, and $BC=2$. The angle bisector of $\angle A$ intersects side $\overline{BC}$ at $D$. What is $BD$?


 
                    \begin{figure}[H]
                        \centering
                        \includegraphics[width=0.5\textwidth]{Problem/AMC_10_10B_2009_20_Problem_Figure_1.png}
                    \end{figure}
                    


 $\text{(A) } \frac {\sqrt3 - 1}{2} \qquad \text{(B) } \frac {\sqrt5 - 1}{2} \qquad \text{(C) } \frac {\sqrt5 + 1}{2} \qquad \text{(D) } \frac {\sqrt6 + \sqrt2}{2} \qquad \text{(E) } 2\sqrt 3 - 1$


 
 \textbf{Answer: }B



\textbf{Solution 1}

By the Pythagorean Theorem, $AC=\sqrt5$. Then, from the Angle Bisector Theorem, we have:

 
 \[\frac{BD}{1}=\frac{2-BD}{\sqrt5}\]
 \[BD\left(1+\frac{1}{\sqrt5}\right)=\frac{2}{\sqrt5}\]
 \[BD(\sqrt5+1)=2\]
 \[BD=\frac{2}{\sqrt5+1}=\boxed{\frac{\sqrt5-1}{2}} \implies \boxed{B}\]

 



\textbf{Solution 2}

Let $\theta = \angle BAD = \angle DAC$. Notice $\tan \theta = BD$ and $\tan 2 \theta = 2$. By the double angle identity, 
 \[2 = \frac{2 BD}{1 - BD^2} \implies BD = \boxed{\frac{\sqrt5 - 1}{2} \implies B}.\]

 Remarks: You could also use tangent half angle formula

 



\textbf{Solution 3}

Let $BD=y$.

 Make $DE$ a line so that it is perpendicular to $AC$. Since $AD$ is an angle bisector, $\triangle AED$ is congruent to $\triangle ABD$.Using the Pythagorean Theorem:

 $AC^2=1^2+2^2$

 $AC^2=5$

 $AC=\sqrt{5}$

 We know that $AE=1$ by the congruent triangles, so $EC=\sqrt{5}-1$. We know that $DE=y$, $EC=\sqrt{5}-1$, and $DC=2-y$. We now have right triangle $DEC$ and its 3 sides. Using the Pythagorean Thereom, we get:

 $y^2+(\sqrt{5}-1)^2=(2-y)^2$

 $-4y=2-2\sqrt{5}$

 So, $y=BD=\boxed{\frac{\sqrt5-1}{2} \implies B}.$

 


	\item (\textit{AMC 10 10B 2009}) What is the remainder when $3^0 + 3^1 + 3^2 + \cdots + 3^{2009}$ is divided by 8?


 $\text{(A) } 0 \qquad \text{(B) } 1 \qquad \text{(C) } 2 \qquad \text{(D) } 4 \qquad \text{(E) } 6$


 
 \textbf{Answer: }D



\textbf{Solution 1}

The sum of any four consecutive powers of 3 is divisible by $3^0 + 3^1 + 3^2 +3^3 = 40$ and hence is divisible by 8. Therefore

 
            \begin{enumerate}[label=(\roman*)]
                \item $(3^2 + 3^3 + 3^4 + 3^5) + \cdots + (3^{2006} + 3^{2007} + 3^{2008} + 3^{2009})$
            \end{enumerate}
            is divisible by 8.  So the required remainder is $3^0 + 3^1 = \boxed {4}$. The answer is $\mathrm{(D)}$.

 



\textbf{Solution 2}

We have $3^2 = 9 \equiv 1 \pmod 8$. Hence for any $k$ we have $3^{2k}\equiv 1^k = 1 \pmod 8$, and then $3^{2k+1} = 3\cdot 3^{2k} \equiv 3\cdot 1 = 3  \pmod 8$.

 Therefore our sum gives the same remainder modulo $8$ as $1 + 3 + 1 + 3 + 1 + \cdots + 1 + 3$. There are $2010$ terms in the sum, hence there are $2010/2 = 1005$ pairs of $1+3$, and thus the sum is $1005 \cdot 4 = 4020 \equiv 20 \equiv \boxed{4} \pmod 8$.

 



\textbf{Solution 3}

We have the formula $\frac{a(r^n-1)}{r-1}$ for the sum of a finite geometric sequence which we want to find the residue modulo 8.
 \[\frac{1 \cdot (3^{2010}-1)}{2}\]
 \[\frac{3^{2010}-1}{2} = \frac{9^{1005}-1}{2}\]
 \[\frac{9^{1005}-1}{2} \equiv \frac{1^{1005}-1}{2} \equiv \frac{0}{2} \pmod 8\]Therefore, the numerator of the fraction is divisible by $8$. However, when we divide the numerator by $2$, we get a remainder of $4$ modulo $8$, giving us $\mathrm{(D)}$.

 Note: you need to prove that 
 \[\frac{9^{1005}-1}{2}\] is not congruent to 0 mod 16 because if so, then the whole thing would be congruent to 0 mod 8, even after dividing by 2 ~ ilikepi12

 



\textbf{Solution 4 (Patterns)}

We can see that $\frac{3^0}{8}$ is has a remainder of 1. $\frac{3^1}{8}$ has a remainder of 4. $\frac{3^2}{8}$ also has a remainder of 4. $\frac{3^3}{8}$ has a remainder of 0.

 If we keep on going, we can see that this pattern repeats every 4 powers of 3 - that is to say, the reminder cycles through $1$, $4$, $4$, and $0$.

 Using simple math, we can see that if assume the pattern repeats every 4 numbers, $3^{2009}$ has a remainder of $\boxed {4}$. The answer is $\mathrm{(D)}$.

 


	\item (\textit{AMC 10 10B 2009}) A cubical cake with edge length $2$ inches is iced on the sides and the top. It is cut vertically into three pieces as shown in this top view, where $M$ is the midpoint of a top edge. The piece whose top is triangle $B$ contains $c$ cubic inches of cake and $s$ square inches of icing. What is $c+s$?


 
                    \begin{figure}[H]
                        \centering
                        \includegraphics[width=0.5\textwidth]{Problem/AMC_10_10B_2009_22_Problem_Figure_1.png}
                    \end{figure}
                    


 $\text{(A) } \frac{24}{5} \qquad \text{(B) } \frac{32}{5} \qquad \text{(C) } 8+\sqrt5 \qquad \text{(D) } 5+\frac{16\sqrt5}{5} \qquad \text{(E) } 10+5\sqrt5$


 
 \textbf{Answer: }B



\textbf{Solution}


                    \begin{figure}[H]
                        \centering
                        \includegraphics[width=0.5\textwidth]{Solution/AMC_10_10B_2009_22_Solution_Figure_1.png}
                    \end{figure}
                    

 Let's label the points as in the picture above. Let $[RNQ]$ be the area of $\triangle RNQ$. Then the volume of the corresponding piece is $c=2[RNQ]$. This cake piece has icing on the top and on the vertical side that contains the edge $QR$. Hence the total area with icing is $[RNQ]+2^2 = [RNQ]+4$. Thus the answer to our problem is $3[RNQ]+4$, and all we have to do now is determine $[RNQ]$.

 



\textbf{Solution 1}

Introduce a coordinate system where $Q=(0,0)$, $P=(2,0)$ and $R=(0,2)$. 

 In this coordinate system we have $M=(2,1)$, and the line $QM$ has the equation $2y-x=0$. 

 As the line $RN$ is orthogonal to $QM$, it must have the equation $y+2x+q=0$ for some suitable constant $q$. As this line contains the point $R=(0,2)$, we have $q=-2$.

 Substituting $x=2y$ into $y+2x-2=0$, we get $y=\frac 25$, and then $x=\frac 45$. 

 We can note that in $\triangle RNQ$ $x$ is the height from $N$ onto $RQ$, hence its area is $[RNQ] = \frac{x \cdot RQ} 2 = \frac{2x}2 = x = \frac 45$, and therefore the answer is $3[RNQ]+4 = 3\cdot \frac 45 + 4 = \boxed{\frac{32}5 \Longrightarrow B}$.

 



\textbf{Solution 2}

Extend $RN$ to intersect $PQ$ at $O$:

 
                    \begin{figure}[H]
                        \centering
                        \includegraphics[width=0.5\textwidth]{Solution/AMC_10_10B_2009_22_Solution_Figure_2.png}
                    \end{figure}
                    

 It is now obvious that $O$ is the midpoint of $PQ$. (Imagine rotating the square $PQRS$ by $90^\circ$ clockwise around its center. This rotation will map the segment $MQ$ to a segment that is orthogonal to $MQ$, contains $R$ and contains the midpoint of $PQ$.

 From $\triangle PQM$ we can compute that $QM = \sqrt{1^2 + 2^2} = \sqrt 5$.

 Observe that $\triangle PQM$ and $\triangle NQO$ have the same angles and therefore they are similar. The ratio of their sides is $\frac{QM}{OQ} = \frac{\sqrt 5}1 = \sqrt 5$.

 Hence we have $ON = \frac{PM}{\sqrt 5} = \frac 1{\sqrt 5}$, and $NQ = \frac{PQ}{\sqrt 5} = \frac 2{\sqrt 5}$.

 Knowing this, we can compute the area of $\triangle NQO$ as $[NQO] = \frac{ON \cdot NQ}2 = \frac 15$.

 Finally, we compute $[RNQ] = [ROQ] - [NQO] = 1 - \frac 15 = \frac 45,$ and conclude that the answer is $3[RNQ]+4 = 3\cdot \frac 45 + 4 = \boxed{\frac{32}5}.$

 
                \begin{itemize}
                    \item You could also notice that the two triangles $\triangle PMQ$ and $\triangle NQR$ in the original figure are similar.
                \end{itemize}


                



\textbf{Solution 3 (Pythagorean Theorem only)}


                    \begin{figure}[H]
                        \centering
                        \includegraphics[width=0.5\textwidth]{Solution/AMC_10_10B_2009_22_Solution_Figure_3.png}
                    \end{figure}
                    

 Since $PQ = SR = 2$ and $PM = MS = 1$, we know that $MQ = MR = \sqrt{2^{2} + 1^{2}} = \sqrt{5}$.  If we let $NQ = x$, then $MN = \sqrt{5} - x$.  Now, by the Pythagorean Theorem, we have:

 
 \[x^{2} + NR^{2} = 2^{2} = 4\]
 \[(\sqrt{5} - x)^{2} + NR^{2} = (\sqrt{5})^{2} = 5\]

 Expanding and rearranging the second equation gives:

 
 \[5 - 2x\sqrt{5} + x^{2} + NR^{2} = 5\]
 \[x^{2} + NR^{2} - 2x\sqrt{5} = 0\]
 \[x^{2} + NR^{2} = 2x\sqrt{5}\]

 Since $x^{2} + NR^{2} = 4$, we have that:

 
 \[2x\sqrt{5} = 4\]
 \[x\sqrt{5} = 2\]
 \[x = \frac{2}{\sqrt{5}} = \frac{2\sqrt{5}}{5}\]

 Knowing $x$, we can solve for the height $NR$:

 
 \[NR^{2} = 2^{2} - x^{2} = 4 - ({\frac{2\sqrt{5}}{5}})^{2} = 4 - \frac{4}{5} = \frac{16}{5}\]
 \[NR = \sqrt{\frac{16}{5}} = \frac{4}{\sqrt{5}} = \frac{4\sqrt{5}}{5}\]

 Therefore, the area of triangle $RNQ$ is $\frac{1}{2} \cdot \frac{2\sqrt{5}}{5} \cdot\frac{4\sqrt{5}}{5} = \frac{1}{2} \cdot\frac{40}{25} = \frac{4}{5}$.  Since the solution to the problem is $3[RNQ] + 4$, the answer is $3(\frac{4}{5}) + 4 = \frac{12}{5} + \frac{20}{5} = \boxed{(B) \frac{32}{5}}$.

 



\textbf{Solution 4}


                    \begin{figure}[H]
                        \centering
                        \includegraphics[width=0.5\textwidth]{Solution/AMC_10_10B_2009_22_Solution_Figure_1.png}
                    \end{figure}
                    

 $MQ = \sqrt{2^{2} + 1^{2}} = \sqrt{5}$

 since $\angle PQM + \angle PMQ = 90 = \angle PQM + \angle NQR$

 therefore $\angle PMQ = \angle NQR$

 and since $\angle MPQ = \angle QNR = 90$

 therefore $\triangle MPQ \sim \triangle QNR$ 

 therefore $\frac {[QNR]}{[MPQ]} = (\frac{QR}{QM})^{2}, [QNR] = [MPQ] \cdot (\frac{QR}{QM})^{2}$

 $[MPQ] = \frac{1}{2} \cdot 2 \cdot 1 = 1$

 $[QNR] = 1 \cdot (\frac{2}{\sqrt{5}})^{2} = 1 \cdot \frac{4}{5} = \frac{4}{5}$

 Since the solution to the problem is $3[QNR] + 4$, the answer is $3(\frac{4}{5}) + 4 = \frac{12}{5} + \frac{20}{5} = \boxed{(B) \frac{32}{5}}$.

 


 



\textbf{Solution 5 (Similarity)}


                    \begin{figure}[H]
                        \centering
                        \includegraphics[width=0.5\textwidth]{Solution/AMC_10_10B_2009_22_Solution_Figure_1.png}
                    \end{figure}
                    

 All units of length in the following solution are in inches, or inches squared, or inches cubed. Units of angles are in degrees.

 
$PQ = 2$. Since $M$ is the midpoint of $\overline{SP}$ which measures $2$, $MP = 1$.

 
Since angle MNR is right, angle QNR is also right. Let $m\angle PQM = x$. Then $m\angle PMQ = 90 - x$. Notice also since $\angle PQR$ is right, $m \angle NQR = 90 - x$. Since $\angle QNR$ is right, $m\angle QRN = x$. Therefore, $\triangle PQM \sim \triangle NRQ$. 

 
Let $QN = a$. By the Pythagorean theorem, $MQ = \sqrt{5}$. By similarity, $\frac{PM}{MQ} = \frac{NQ}{QR} \longrightarrow \frac{1}{\sqrt{5}} = \frac{a}{2}$, so $a=\frac{2}{\sqrt{5}}=\frac{2\sqrt{5}}{5}$. By the Pythagorean theorem, $NR^2+NQ^2=QR^2$. Substituting known values in and solving for $NR$, we get $NR=\frac{4\sqrt{5}}{5}$. (Alternatively, use the fact that $\triangle PQM \sim \triangle NRQ$). Since $\triangle NQR$ is a right triangle, the area is just $NQ \cdot NR \cdot \frac{1}{2}$ which, substituting values, is equal to $\frac{4}{5}$. But remember that $s$ also consists of the side of the cake, so we have to add $2^2=4$. So $s=\frac{4}{5}+4=\frac{24}{5}$. 

 
Meanwhile, $c$ is the volume of the slice (a triangular prism) which is found by the base area times height. We already calculated the base area to be $\frac{4}{5}$, so simply multiply by $2$ to get the volume $=\frac{8}{5}$. This is the value of $c$.

 Sum $c+s$: $c+s=\frac{8}{5}+\frac{24}{5}=\frac{32}{5} \Longrightarrow \boxed{\textbf{(B) } \frac{32}{5}}$.

 



\textbf{Solution 6 (only Pythagorean Theorem, no algebra)}

Label the vertices of the square, $P$, $Q$, $R$, and $S$ and draw line segment $MA$ (as shown below):
                    \begin{figure}[H]
                        \centering
                        \includegraphics[width=0.5\textwidth]{Solution/AMC_10_10B_2009_22_Solution_Figure_6.png}
                    \end{figure}
                    

 $PM=MS=1$ and $SR=2$, so by the pythagorean theorem, 
 \[MR = \sqrt{1^2+2^2} = \sqrt{5}.\] By the same logic, $MQ=\sqrt{5}$. The area of $\triangle QMR$ is the area of the whole square, minus the combined areas of $\triangle PMQ$ and $\triangle MSR$, so 
 \[[QMR] = 4-1-1=2.\] Since $QM$ is the base of $\triangle QMR$ and $RN$ is the height, 
 \[\frac{QM(RN)}{2} = 2\] 
 \[\frac{\sqrt5(RN)}{2} = 2,\] so $RN=\frac{4\sqrt5}{5}$. We also know that 
 \[QN = QM-MN = \sqrt{5} - \frac{3\sqrt5}{5} = \frac{2\sqrt5}{5}.\] Now, we can find the area of $\triangle QNR$. 
 \[[QNR] = \frac{1}{2}(QN)(RN) = \frac{1}{2}(\frac{4\sqrt5}{5})(\frac{2\sqrt5}{5}) = \frac{4}{5}.\] The area of icing is the area of $\triangle QNR$, plus the area of the 2x2 square on $QR$, so 
 \[s=\frac{4}{5}+4 =\frac{4}{5} + \frac{20}{5} = \frac{24}{5}.\] The cubic inches of cake is the volume of the piece, which is the area of $\triangle QNR$ times $2$ (the height of the cake), so $c = \frac{8}{5}$. Hence, $c+s = \frac{32}{5}$, and the answer is $\boxed{\textbf{(B) } \frac{32}{5}}$

 



\textbf{Solution 7 (Geometric Shifting & The Right Triangle Altitude Theorem)}

Extend line segment $QM$ past $M$ until it intersects the line containing side $RS$ at point $T$:

 
                    \begin{figure}[H]
                        \centering
                        \includegraphics[width=0.5\textwidth]{Solution/AMC_10_10B_2009_22_Solution_Figure_7.png}
                    \end{figure}
                    

 Consider $\triangle PMQ$ and $\triangle SMT$:

 
                \begin{itemize}
                    \item $\angle MPQ = \angle MST = 90^\circ$ (interior angles of a square and its linear extension).\item $PM = MS = 1$ (since $M$ is the midpoint of $PS$).\item $\angle PMQ = \angle SMT$ (vertical angles).
                \end{itemize}


                By Angle-Side-Angle (ASA Congruence), $\triangle PMQ \cong \triangle SMT$. By Corresponding Parts of Congruent Triangles (CPCTC), we have $ST = PQ = 2$. Therefore, the total base length of the extended structure is:
 \[RT = RS + ST = 2 + 2 = 4\]

 Now examine the large right triangle $\triangle RTQ$, tracking angles from the right angle ($R$), to the shallow angle ($T$), to the sharp angle ($Q$):

 
                \begin{itemize}
                    \item $\angle QRT = 90^\circ$ because the horizontal base extension line $RT$ meets the vertical square side $QR$ at a right angle.\item $RN$ is an altitude drawn from the right angle $R$ to the hypotenuse $QT$, since $RN \perp MQ$ and points $Q, N, M, T$ are collinear.
                \end{itemize}


                By the Right Triangle Altitude Theorem, matching vertices strictly by their corresponding angle measures ($\text{Right Angle} \rightarrow \text{Angle } T \rightarrow \text{Remaining Angle}$):

 
            \begin{enumerate}
                \item $\triangle NTR \sim \triangle RTQ$ (sharing acute angle $\angle NTR$ and a right angle).\item $\triangle NRQ \sim \triangle RTQ$ (sharing acute angle $\angle RQT$ and a right angle).\item By the Law of Syllogism (Transitivity of Similarity), $\triangle NTR \sim \triangle NRQ$.
            \end{enumerate}
            
Using the Pythagorean Theorem on the master triangle $\triangle RTQ$:
 \[QT = \sqrt{QR^2 + RT^2} = \sqrt{2^2 + 4^2} = \sqrt{20} = 2\sqrt{5}\]

 Because $\triangle NRQ \sim \triangle RTQ$, the ratio of the \textbf{Hypotenuse} to the \textbf{Long Leg} must be identical in both triangles:
 \[\frac{QR}{NR} = \frac{QT}{RT} \implies \frac{2}{NR} = \frac{2\sqrt{5}}{4}\]

 Flipping the proportion yields:
 \[\frac{NR}{2} = \frac{4}{2\sqrt{5}} \implies NR = \frac{4}{\sqrt{5}}\]

 Applying the Pythagorean Theorem on $\triangle NRQ$ to isolate the short leg $NQ$:
 \[NQ = \sqrt{QR^2 - NR^2} = \sqrt{2^2 - \left(\frac{4}{\sqrt{5}}\right)^2} = \sqrt{4 - \frac{16}{5}} = \sqrt{\frac{4}{5}} = \frac{2}{\sqrt{5}}\]

 We can now evaluate the area of triangle $B$ ($\triangle NRQ$):
 \[[NRQ] = \frac{1}{2} \cdot NR \cdot NQ = \frac{1}{2} \cdot \frac{4}{\sqrt{5}} \cdot \frac{2}{\sqrt{5}} = \frac{4}{5}\]

 Since the cake has a height of $2$ inches, the piece contains $c = 2[NRQ] = \frac{8}{5}$ cubic inches of cake. The icing covers the top face and the outer $2 \times 2$ vertical wall, giving $s = [NRQ] + 4 = \frac{24}{5}$ square inches of icing. 

 Evaluating the final target sum $c + s$:
 \[3\left(\frac{4}{5}\right) + 4 = \frac{12}{5} + \frac{20}{5} = \boxed{\text{(B) } \frac{32}{5}}\]

 


 


	\item (\textit{AMC 10 10B 2009}) Rachel and Robert run on a circular track.  Rachel runs counterclockwise and completes a lap every 90 seconds, and Robert runs clockwise and completes a lap every 80 seconds.  Both start from the same line at the same time.  At some random time between 10 minutes and 11 minutes after they begin to run, a photographer standing inside the track takes a picture that shows one-fourth of the track, centered on the starting line.  What is the probability that both Rachel and Robert are in the picture?


 $\text{(A) } \frac 1{16} \qquad \text{(B) } \frac 18 \qquad \text{(C) } \frac 3{16} \qquad \text{(D) } \frac 14 \qquad \text{(E) } \frac 5{16}$


 
 \textbf{Answer: }C



\textbf{Solution}

After $10$ minutes $(600$ seconds$),$ Rachel will have completed $6$ laps and be $30$ seconds from completing her seventh lap. Because Rachel runs one-fourth of a lap in $22.5$ seconds, she will be in the picture between $18.75$ seconds and $41.25$ seconds of the tenth minute.  After 10 minutes, Robert will have completed $7$ laps and be $40$ seconds from completing his eighth lap.  Because Robert runs one-fourth of a lap in $20$ seconds, he will be in the picture between $30$ seconds and $50$ seconds of the tenth minute.  Hence both Rachel and Robert will be in the picture if it is taken between $30$ seconds and $41.25$ seconds of the tenth minute.  So the probability that both runners are in the picture is $\frac {41.25 - 30} {60} = \boxed{\frac {3}{16}}$.  The answer is $\mathrm{(C)}$.

 


	\item (\textit{AMC 10 10B 2009}) The keystone arch is an ancient architectural feature. It is composed of congruent isosceles trapezoids fitted together along the non-parallel sides, as shown. The bottom sides of the two end trapezoids are horizontal. In an arch made with $9$ trapezoids, let $x$ be the angle measure in degrees of the larger interior angle of the trapezoid. What is $x$? 


 
                    \begin{figure}[H]
                        \centering
                        \includegraphics[width=0.5\textwidth]{Problem/AMC_10_10B_2009_24_Problem_Figure_1.png}
                    \end{figure}
                    


 $\text{(A) } 100 \qquad \text{(B) } 102 \qquad \text{(C) } 104 \qquad \text{(D) } 106 \qquad \text{(E) } 108$


 
 \textbf{Answer: }A



\textbf{Solution 1}

Extend all the legs of the trapezoids. They will all intersect in the middle of the bottom side of the picture, forming the situation shown below.

 
                    \begin{figure}[H]
                        \centering
                        \includegraphics[width=0.5\textwidth]{Solution/AMC_10_10B_2009_24_Solution_Figure_1.png}
                    \end{figure}
                    

 Each of the angles at $X$ is $\frac{180^\circ}9 = 20^\circ$. From $\triangle XYZ$, the degree measure of the smaller interior angle of the trapezoid is $\frac{180^\circ - 20^\circ}2 = 80^\circ$, hence the degree measure of the larger interior angle is $180^\circ - 80^\circ = \boxed{100^\circ \Longrightarrow A}$.

 Proof that all the extended trapezoid legs intersect at the same point: It is sufficient to prove this for any pair of neighboring trapezoids. For two neighboring trapezoids, the situation is symmetric according to their common leg, therefore the extensions of both outside legs intersect the extension of the common leg at the same point, Q.E.D.

 Knowing this, we can now easily see that the intersection point must be on the bottom side of our picture, as it lies on the bottom leg of the rightmost trapezoid. And by symmetry the point must be in the center of this side.

 



\textbf{Solution 2}

A decagon can be formed from the trapezoids and the base. The sum of the decagon's angles is $180(10-2)=1440^\circ$. Letting the larger angle in each trapezoid be $x$, the two angles formed by the line each measures $(180-x)^\circ$. There are $8$ congruent angles left. Each of those angles measures $(360-2x)^\circ$.  Putting it all together: $8(360-2x)+2(180-x)=1440\implies x=\boxed{\mathrm{(A)\ }100^\circ}$.

 



\textbf{Solution 3}

If we reflect the arch across the line, we form an 18-gon. $\frac{180*(18-2)}{18} = 160^\circ$ so each interior angle of the 18-gon is $160^\circ$. Let $x$ be the degree measure of the larger interior angle of a trapezoid. From the diagram, we see that $2x + 160 = 360$, so $2x = 200$ and $x = 100$, or $\boxed{A}$.

 


	\item (\textit{AMC 10 10B 2009}) Each face of a cube is given a single narrow stripe painted from the center of one edge to the center of the opposite edge.  The choice of the edge pairing is made at random and independently for each face.  What is the probability that there is a continuous stripe encircling the cube?


 $\text{(A) } \frac 18 \qquad \text{(B) } \frac {3}{16} \qquad \text{(C) } \frac 14 \qquad \text{(D) } \frac 38 \qquad \text{(E) } \frac 12$


 
 \textbf{Answer: }B



\textbf{Solution 1}

There are two possible stripe orientations for each of the six faces of the cube, so there are $2^6 = 64$ possible stripe combinations.  There are three pairs of parallel faces so, if there is an encircling stripe, then the pair of faces that does not contribute uniquely determines the stripe orientation for the remaining faces.  In addition, the stripe on each face that does not contribute may be oriented in either of two different ways. So a total of $3 \cdot 2 \cdot 2 = 12$ stripe combinations on the cube result in a continuous stripe around the cube.  The required probability is $\frac {12}{64} = \boxed{\frac {3}{16}}$.

 



\textbf{Solution 2}

Without loss of generality, orient the cube so that the stripe on the top face goes from front to back.  There are two mutually exclusive ways for there to be an encircling stripe: either the front, bottom and back faces are painted to complete an encircling stripe with the top face's stripe or the front, right, back and left faces are painted to form an encircling stripe.  The probability of the first case is $\left(\frac 12\right)^3 = \frac 18$, and the probability of the second case is $\left(\frac 12\right)^4 = \frac {1}{16}$.  The cases are disjoint, so the probabilities sum to $\frac 18 + \frac {1}{16} = \boxed {\frac {3}{16}}$.

 



\textbf{Solution 3}

There are three possible orientations of an encircling stripe (as each encircling would leave out 1 pair of opposite faces, and there are 3 pairs in total).  For any one of these to appear, the stripes on the four faces through which the continuous stripe is to pass must be properly aligned.  The probability of each such stripe alignment is $\left(\frac 12\right)^4 = \frac {1}{16}$.  Since there are three such possibilities and they are disjoint, the total probability is $3 \cdot \frac {1}{16} = \boxed{\frac {3}{16}}$.  The answer is $\mathrm{(B)}$.

 



\textbf{Solution 4}

Consider a vertex on the cube and the three faces that are adjacent to that vertex. If no two stripes on those three faces are aligned, then there is no stripe encircling the cube. The probability that the stripes aren't aligned is $\frac{1}{4}$, since for each alignment of one stripe, there is one and only one way to align the other two stripes out of four total possibilities. therefore there is a probability of $\frac{3}{4}$ that two stripes are aligned.

 Now consider the opposing vertex and the three sides adjacent to it. Given the two connected stripes next to our first vertex, we have two more that must be connected to make a continuous stripe. There is a probability of $\left(\frac{1}{2}\right)^2=\frac{1}{4}$ that they are aligned, so there is a probability of $\frac{3}{4}\cdot \frac{1}{4}=\boxed{\frac {3}{16}}$ that there is a continuous stripe.

 


 



\textbf{Solution 5}

Each face has two possible stripes, so there are $2^6=64$ ways in total.

 We need to choose $2$ opposite faces to leave out. Since the second is fixed once we pick the first face, we only need to calculate $\binom{3}{1}=3$ orientations. Next, the two faces also each have 2 orientations, since they don't matter to the stripe. Therefore there are $3\cdot2\cdot2=12$ ways. The probability is $\frac{12}{64}=\frac{3}{16} \Longrightarrow \boxed{\textbf{(B) } \frac{3}{16}}$.

 


\end{enumerate}